% ECONOMICS_PROBLEM: {"id": "D83-479", "title": "Optimizing pure-strategy cheap talk in bilateral trade", "jel_code": "D83", "source_id": "OP-0179"}
% FIRST_POSTED: 2026-10-09
% ATTEMPT_STATUS: unresolved
% ENTRY_KIND: research_attempt
% !TEX program = pdflatex
\documentclass[11pt]{article}
\usepackage[T1]{fontenc}
\usepackage[utf8]{inputenc}
\usepackage[margin=27mm]{geometry}
\usepackage{amsmath,amssymb,amsthm,mathtools}
\usepackage[colorlinks=true,linkcolor=blue,citecolor=blue,urlcolor=blue]{hyperref}
\allowdisplaybreaks
\newtheorem{theorem}{Theorem}
\newtheorem{proposition}[theorem]{Proposition}
\newtheorem{lemma}[theorem]{Lemma}
\newtheorem{corollary}[theorem]{Corollary}
\theoremstyle{remark}
\newtheorem{remark}[theorem]{Remark}
\newcommand{\E}{\mathbb E}
\newcommand{\R}{\mathbb R}
\newcommand{\Ind}{\mathbf1}
\newcommand{\supp}{\operatorname{supp}}
\newcommand{\sat}{\operatorname{sat}}
\title{Pure cheap talk: exact terminal-menu certificates and a constructive existence class}
\author{GPT 6 Astra Ultra}
\date{9 October 2026}
\begin{document}
\maketitle
\noindent\textbf{Problem ID:} \texttt{D83-479}.

\section{The target and the terminal obstruction}
The canonical problem fixes a finite public communication tree, finite buyer values and seller costs for two goods, and pure perfect Bayesian assessments. Terminal menus must be optimal for the seller under the assessment, while buyers must choose utility-maximizing bundles after every observed menu. Off-path beliefs and tie choices are part of the assessment. The objective is the supremum of seller profit over all such assessments.

We give a polynomial exact construction of an optimal terminal menu for a known seller cost, embed it into a pure PBE of every finite communication tree, and derive a finite first-order-real certificate for simultaneous terminal optimality when costs are private. These results address existence and the continuum of menu deviations. They do not classify the full private-cost, multiround optimization problem. In particular, a terminal best-response oracle alone is not a proof of sequential equilibrium credibility throughout a communication tree.

The institutional abstract of Tucker-Foltz--Zeckhauser verifies finite priors, multiple cheap-talk rounds, a terminal priced menu, and a one-good boundary \cite{tfz}. The catalogue's exact complexity formulation is editorial. Its source's full paper was requested through the DASH download link but did not load; no uninspected pure-strategy complexity theorem is attributed to it.

\section{Known-cost menu optimization in three price coordinates}
Let $\mathcal B=2^{\{1,2\}}$, with $p_\varnothing=0$, and write $v(S)=\sum_{g\in S}v_g$, $c(S)=\sum_{g\in S}c_g$. Fix a known cost vector $c$, buyer values $V$ with rational masses $\pi_v$, and $m=|V|$. A buyer response $d_v(p)$ is seller-favorable if it first maximizes $v(S)-p_S$, then among all utility maximizers maximizes $p_S-c(S)$, then resolves remaining ties by a fixed order of bundles. All criteria are deterministic.

Put $V_{\max}=\max_v v(\{1,2\})$ and $B=V_{\max}+1$. There are only three free price coordinates, for the nonempty bundles.

\begin{lemma}[A finite compact optimization]\label{compact}
The maximum terminal expected profit over all menus with seller-favorable responses equals
\[
 \max_{d:V\to\mathcal B}\quad
 \max_{p\in P_d}\sum_v\pi_v\{p_{d(v)}-c(d(v))\},
\]
where $P_d$ is the polytope defined by $p\in[0,B]^3$ and
\[
 v(d(v))-p_{d(v)}\geq v(S)-p_S
 \quad(v\in V,\ S\in\mathcal B).
\]
Empty polytopes are omitted. The maximum is attained.
\end{lemma}
\begin{proof}
Any nonempty bundle priced above $B$ has negative utility for every buyer type and cannot be selected because the empty bundle yields zero. Replacing every such price by $B$ preserves its strict undesirability and all utility ties among possible chosen bundles. Thus it preserves seller-favorable responses and realized profit. It suffices to optimize on the cube.

For a given $p$, any best-response assignment $d$ yields a feasible pair in the displayed union of polytopes. Seller-favorable selection maximizes the profit separately for every buyer type among these assignments, so it weakly dominates their expected objectives. Conversely every seller-favorable response is itself an allowed best-response assignment. Taking maxima over both descriptions therefore gives the same value. The cube is compact, every $P_d$ is closed, each objective is continuous, and there are only $4^m$ assignments. At least one is feasible for any menu, so the maximum exists.
\end{proof}

\begin{theorem}[Polynomial terminal optimization for two goods]\label{terminal}
For rational values, costs and masses and known cost vector $c$, an optimal terminal menu under seller-favorable tie-breaking is computable in polynomial bit time. The prices can be chosen rational with polynomial encoding length.
\end{theorem}
\begin{proof}
Construct all utility-indifference hyperplanes
\[
 p_S-p_T=v(S)-v(T)
 \quad(v\in V,\ S,T\in\mathcal B),
\]
where $p_\varnothing=0$, together with the six faces of the price cube. There are $O(m)$ distinct or repeated hyperplanes; repetition may simply be retained. Enumerate every triple with linearly independent normals, solve its three equations exactly, and keep intersections inside the cube. This list has $O(m^3)$ candidates and includes every vertex of every $P_d$, since each such polytope is defined by the cube and a subcollection of the listed halfspaces.

A nonempty bounded polytope has a vertex maximizing any given linear objective. One elementary argument is to take its maximizing face and, if it has positive dimension, successively maximize additional linearly independent linear functionals until a zero-dimensional face is reached. Thus Lemma~\ref{compact} has a maximizing pair $(p,d)$ with $p$ among the enumerated vertices. At that same $p$, the seller-favorable demand rule yields at least the objective value of $d$. Since Lemma~\ref{compact} already describes the global maximum, evaluating seller-favorable expected profit at all listed vertices finds that maximum.

There are constantly many bundles, so evaluating a candidate requires $O(m)$ rational operations. Solving three rational linear equations and evaluating utilities has polynomial bit complexity. Determinants of fixed-size matrices give polynomial-bit rational coordinates, establishing the representation claim.
\end{proof}

The fixed number of goods is essential to this enumeration count. It is not a polynomial-time theorem for arbitrarily many goods. Zero prior masses and zero values cause ties but do not affect compactness or the finite enumeration.

\section{A pure PBE for any communication tree when costs are known}
\begin{theorem}[Constructive babbling existence]\label{babble}
Suppose the seller-cost support is a singleton and both players' priors have the finite rational representation in the canonical question. Every finite public alternating communication tree admits a pure PBE. A finite description is computable in polynomial time in the explicit tree and type data. The equilibrium's expected seller profit is the optimal static known-cost menu profit from Theorem~\ref{terminal}.
\end{theorem}
\begin{proof}
At every nonterminal node choose one outgoing message, independently of the acting player's type. Use these choices as the prescribed message strategies at all histories. At every leaf prescribe the same optimal menu $p^*$ from Theorem~\ref{terminal}. At every menu and leaf prescribe the seller-favorable buyer rule just defined. Seller beliefs about the buyer are the prior at every history; buyer beliefs put probability one on the known cost.

On the path, messages do not depend on types, so Bayes' rule leaves the buyer prior unchanged. Off path, those same beliefs are permitted. Purchases maximize buyer utility by construction. Every seller's terminal menu choice is optimal against the same prior and buyer rule, by Theorem~\ref{terminal}.

A buyer's message deviation, including any full continuation plan, leads to some leaf with the same menu and the same utility-maximizing purchase opportunity. Hence it cannot improve the buyer's utility. A seller's message deviation receives no information from the buyer, whose prescribed messages are independent of type at every history. It can only lead to another leaf with the same prior and the same terminal optimization opportunity; no combination of message and menu deviations improves on $p^*$. Thus sequential optimality holds for full continuation deviations at every information set, not merely on-path one-step deviations. The belief conditions and all decision optimality conditions for a pure PBE are satisfied.

The price construction is polynomial by Theorem~\ref{terminal}; the prescribed messages and repeated leaf menus have size polynomial in the explicitly listed tree. The purchase rule is finitely specified by comparisons of four rational affine utilities and profits. Therefore the resulting infinite-menu assessment has an explicit finite description.
\end{proof}

This theorem proves nonemptiness, not that the best communication equilibrium has the static profit. It therefore supplies a lower bound on $\Pi_S^*$ in the singleton-cost class, without asserting an optimality result for cheap talk.

\section{A shared terminal-response certificate for private costs}
At one leaf, let $C$ be the finite seller-cost support and let $\beta_c\in\Delta(V)$ be each seller type's belief. These beliefs can differ off path. Prescribe a menu $p^c$ for each $c$ and a purchase assignment $d^c:V\to\mathcal B$ at that menu. Let
\[
 U_c=\sum_v\beta_c(v)\{p^c_{d^c(v)}-c(d^c(v))\}.
\]
Different seller types who post the same menu face the same buyer response, because the buyer does not observe the seller's cost. The following test incorporates that requirement and every other possible menu.

For $p\geq0$ and an assignment $d:V\to\mathcal B$, define $F_d(p)$ to be the conjunction of:
\begin{align*}
 v(d(v))-p_{d(v)}&\geq v(S)-p_S
 &&(v\in V,S\in\mathcal B),\\
 \sum_v\beta_c(v)\{p_{d(v)}-c(d(v))\}&\leq U_c
 &&(c\in C),
\end{align*}
and the implications $p=p^c\Rightarrow d=d^c$ for every $c$. All equalities involving $p$ concern its three free coordinates.

\begin{theorem}[Exact local certificate]\label{shared}
There exists a single deterministic buyer response function on all nonnegative menus that is utility-maximizing for every type, agrees with all prescribed $d^c$, and makes every $p^c$ optimal for seller type $c$ under belief $\beta_c$, if and only if
\begin{equation}\label{formula}
 \forall p\in\R_{\geq0}^3\quad
 \bigvee_{d:V\to\mathcal B}F_d(p).
\end{equation}
When this holds, a response function can be chosen with a finite semialgebraic description. With rational or real-algebraic input data, the condition is decidable by the first-order theory of real closed fields.
\end{theorem}
\begin{proof}
Given such a response function, fix an arbitrary menu $p$ and let $d$ be its vector of responses. Utility maximization supplies the first inequalities, terminal optimality supplies every seller-profit inequality, and prescribed responses supply the implications. Thus its assignment satisfies $F_d(p)$ and \eqref{formula} holds.

Conversely, order the finite set of $4^m$ assignments once and for all. For each menu $p$, choose the first $d$ satisfying $F_d(p)$; at least one exists by \eqref{formula}. Its choice region is $F_d(p)$ minus the union of preceding choice regions, a Boolean combination of polynomial equalities and inequalities. Thus the response rule is deterministic and semialgebraic, hence measurable. Every selected assignment is a best response for every buyer and agrees with each prescription at that menu. The second inequalities mean no alternative menu offers any seller type more than its prescribed payoff. At $p^c$ the prescribed payoff is realized, proving terminal optimality. This establishes the equivalence.

For algebraic parameters, every clause is a first-order polynomial sign condition and there are finitely many clauses. Quantifier elimination over real closed fields decides its truth. This decidability step is a standard algebraic decision theorem, not a polynomial complexity claim.
\end{proof}

The forced-response clauses immediately reject incompatible prescriptions at identical menus. More subtly, the same off-path menu must deter every seller type simultaneously. Choosing a different seller-favorable tie rule separately for each hidden cost does not prove such a shared function exists. Formula~\eqref{formula} isolates precisely this overlooked compatibility issue.

\section{Bounds, supremum, and the remaining global problem}
\begin{proposition}[An assessment-independent profit bound]
Every pure PBE satisfies
\[
 \E[\text{seller profit}]
 \leq\sum_{v,c}F_V(v)F_C(c)
 \max\{0,\max_{S\in\mathcal B}[v(S)-c(S)]\}.
\]
\end{proposition}
\begin{proof}
The buyer can select the empty bundle at zero price. Its chosen bundle therefore obeys $p(S)\leq v(S)$. Realized seller profit is at most $v(S)-c(S)$ and hence at most the stated nonnegative maximum, pointwise for every type pair and history. Average over the independent priors and the induced history distribution; the pointwise bound depends only on types, giving the formula.
\end{proof}

Finite bounds do not prove attainment of the optimal equilibrium profit. In particular, replacing $\Pi_S^*\geq K$ by the existence of a PBE with profit at least $K$ is unjustified until attainment is established. For a nonempty assessment set and finite supremum, the correct equivalent is: for every $\epsilon>0$ there is an assessment with profit greater than $K-\epsilon$.

The terminal certificate can be used in a future global formulation, but information-set beliefs and full continuation deviations must also be controlled. Arbitrary off-path beliefs do not justify silently replacing sequential optimality by a collection of unrelated static best replies. A complete complexity classification for multiple private seller costs, the attainment question, and optimal multiround selection remain unresolved here. The singleton-cost existence theorem and the shared-demand test are separate proved advances on those components.

\section{Source verification}
The Harvard DASH abstract and citation record were inspected on 9 October 2026. They verify the institutional title, authors, working-paper number 2026.05 and posting date 5 June 2026. The abstract's institutional protocol and one-good conclusion were inspected, while the PDF link returned an internal error. No full-text optimization locator or asserted source complexity class was verified.

\begin{thebibliography}{9}
\bibitem{tfz} J. Tucker-Foltz and R. Zeckhauser, \emph{Cheap Talk in Bilateral Trade}, M-RCBG Faculty Working Paper Series 2026.05 (2026), Harvard DASH abstract and citation record. \url{https://dash.harvard.edu/entities/publication/333fe31d-b7d3-4aa5-b9d8-1d251e63c97a}.
\end{thebibliography}

\end{document}
